\Qhead{6.3}{Линейная фильтрация Колмогорова--Винера. Согласованный фильтр}{Линейная фильтрация Колмогорова-Винера на основе минимизации дисперсии ошибки. Принцип ортогональности ошибки и наблюдения. Реализуемые линейные фильтры и уравнение Винера-Хопфа. Выделение сигнала из шума. Согласованный фильтр.}
\Subfile{6.3.1}{sections/06_detection/q03/s01}
\Subfile{6.3.2}{sections/06_detection/q03/s02}
\Subfile{6.3.3}{sections/06_detection/q03/s03}

\recap{Винер минимизирует среднюю ошибку случайного сообщения через ортогональность. Причинность ведёт к Винеру–Хопфу. Согласованный фильтр максимизирует SNR известного сигнала и в гауссовском шуме реализует статистику обнаружения.}
\board{Ортогональность, $H_{\rm nc}$, причинное интегральное уравнение, $H_{\rm matched}$ и $2E/N_0$.}
\oral{Почему винеровский фильтр $S_x/(S_x+S_n)$ и согласованный $S^*/S_n$ решают разные задачи?}
\bait{Можно оставить причинность: обрезание некаузального оптимума обычно не даёт причинного оптимального фильтра.}
\sources{Биленко и соавт.~\cite{bilenko}, §5.8; MIT 6.432~\cite{wornell}, обнаружение и винеровская фильтрация; Wunsch~\cite{wunsch}, §6.1. Выводы ортогональности и SNR приведены полностью в принятой нормировке преобразования Фурье.}
